% Lösungen Einheit 1 von 4 – Das Skalarprodukt als Objekt (zusammenbau v0.1)
\einheitenkopf{Einheit 1 von 4 · Das Skalarprodukt als Objekt}

\erg{9}{a) Zahl \quad b) $\vec a \circ \vec b = 3 > 0$: spitzer Winkel \quad c) Zahl \quad d) $r \le -2$ (aus $\vec u \circ \vec v = r + 2 \le 0$) \quad e) falsch: gleiche Richtung heißt, ein Vektor ist ein positives Vielfaches des anderen; bei gleicher Koordinatensumme ist der Faktor eins, die Vektoren sind also gleich \quad f) Ausmultipliziert bleibt nur $\vec u \circ \vec u = s^2$, denn $\vec u \circ \vec v$ und $\vec u \circ \vec w$ sind null (senkrechte Kanten); h kommt nicht vor}
\erg{10}{a) Fehler: $\vec b \circ \vec c$ ist eine Zahl, und Vektor plus Zahl ist nicht erklärt; richtig ist „nicht definiert“ \quad b) Es ist die Summe der Produkte entsprechender Koordinaten, $a_1 b_1 + a_2 b_2 + a_3 b_3$; Produkte und Summen von Zahlen sind Zahlen}
